\begin{lemma}
\label{lem:p071-exact-star-polyhedral-envelope}
Let \(A\) be finite and \(K\) countable.  Let \(D\) be the compact exact-star
core having exceptional branches \([-R_a,0]\), \(a\in A\), full branches
\([-\infty,0]\), \(b\in K\), and common node \(\eta\).  Let \(Q\) be a
nonempty compact convex polytope.

Write \(\mathcal A(D\times Q)\) for the continuous real-valued functions
whose restrictions to every branch prism are jointly convex and which
satisfy the exact reciprocal-weight product junction inequalities: for every
finite \(F\subset K\), admissible \(r_b>0\), \(0<\rho_a\leq R_a\), and
points \(y_0,y_b,y_a\in Q\) with
\[
 \Lambda y_0=
   \sum_{b\in F}\frac{y_b}{r_b}
   +\sum_{a\in A}\frac{y_a}{\rho_a},
 \qquad
 \Lambda=\sum_{b\in F}\frac1{r_b}
            +\sum_{a\in A}\frac1{\rho_a}>0,
\]
one has
\[
 \Lambda v(\eta,y_0)\leq
   \sum_{b\in F}\frac{v_b(-r_b,y_b)}{r_b}
   +\sum_{a\in A}\frac{v_a(-\rho_a,y_a)}{\rho_a}.
\]
For \(\phi\in C(D\times Q)\), define
\begin{equation}
 C^{\mathrm{ex}}\phi(x)=
 \sup\{v(x):v\in\mathcal A(D\times Q),\ v\leq\phi\}.
 \label{eq:p071-exact-envelope-definition}
\end{equation}
Then \(C^{\mathrm{ex}}\phi\) belongs to
\(\mathcal A(D\times Q)\), is the largest member of that class below
\(\phi\), and satisfies
\begin{equation}
 \lVert C^{\mathrm{ex}}\phi-C^{\mathrm{ex}}\psi\rVert_{D\times Q}
 \leq \lVert\phi-\psi\rVert_{D\times Q}
 \qquad(\phi,\psi\in C(D\times Q)).
 \label{eq:p071-exact-envelope-nonexpanding}
\end{equation}
In particular the envelope is continuous simultaneously at the star node,
the full endpoints, and every face of \(Q\).
\end{lemma}

\begin{proof}
The defining family is nonempty because it contains every constant below
\(\min\phi\).  Hence \(C:=C^{\mathrm{ex}}\phi\) is finite,
\(\min\phi\leq C\leq\phi\), and is lower semicontinuous as a supremum of
continuous functions.  A pointwise supremum preserves joint convexity on
every branch prism.  It also preserves the junction inequalities: apply the
inequality to one competitor before taking its supremum at the centre, and
then bound the supremum of the weighted sum by the weighted sum of the
pointwise suprema.  Since the set of full branches is countable, the common
countable-activity condition is automatic on this reduced core.

We prove the full-endpoint trace together with the uniform estimate needed
for continuity.  Fix \(b\in K\) and \(\delta>0\).  Continuity of \(\phi_b\)
on the compact product \([-\infty,0]\times Q\) gives a \(T<0\) such that
\begin{equation}
 \lvert\phi_b(t,y)-\phi_b(-\infty,y)\rvert<\delta/3
 \quad(t\leq T,\ y\in Q).
 \label{eq:p071-exact-envelope-obstacle-terminal-control}
\end{equation}
Every competitor \(v\) is nondecreasing along a full branch: a finite convex
function of \(t\in(-\infty,0]\) having a finite limit at \(-\infty\) has
nonnegative secant slopes.  Put
\[
                         e_v(y)=v_b(T,y)-\delta
\]
and define \(w\) by \(w=v-\delta\) off branch \(b\), while on that branch
set
\begin{equation}
                  w_b(t,y)=
          \max\{v_b(t,y)-\delta,e_v(y)\}.
 \label{eq:p071-exact-envelope-terminal-repair}
\end{equation}
The function \(e_v\) is continuous and convex.  Moreover
\(e_v\leq v(\eta,\cdot)-\delta\), so the common node value of \(w\) is
still \(v(\eta,\cdot)-\delta\).  Thus \(w\) is continuous and jointly
convex on every branch.  A common downward translation preserves every
junction inequality, and the subsequent increase on one full branch only
increases its right-hand side.  Hence \(w\) satisfies all junction
inequalities.

For \(t\geq T\), monotonicity gives
\(w_b(t,y)=v_b(t,y)-\delta\leq\phi_b(t,y)\).  For \(t\leq T\),
\eqref{eq:p071-exact-envelope-obstacle-terminal-control} gives
\[
 e_v(y)\leq\phi_b(T,y)-\delta
 <\phi_b(-\infty,y)-2\delta/3
 <\phi_b(t,y),
\]
while the other entry in
\eqref{eq:p071-exact-envelope-terminal-repair} is also below \(\phi_b\).
Consequently \(w\leq\phi\) on the whole compact product, and \(w\) is an
admissible competitor.  Its endpoint value is \(e_v\), so
\[
 C(-\infty,y)\geq v_b(T,y)-\delta.
\]
Taking the supremum over \(v\), and using monotonicity once more, yields
\begin{equation}
 0\leq C_b(t,y)-C_b(-\infty,y)\leq\delta
 \quad(t\leq T,\ y\in Q).
 \label{eq:p071-exact-envelope-uniform-terminal-trace}
\end{equation}
Thus \(C\) has the exact endpoint trace, uniformly in the polyhedral
variable.

On every finite branch prism, \(C\) is finite, convex, and lower
semicontinuous.  A finite convex function on a compact polytope is upper
semicontinuous: triangulate locally with the chosen point as a vertex and
let its barycentric coefficient tend to one.  Hence \(C\) is continuous on
every finite branch prism.  The endpoint slice is likewise finite, convex,
and lower semicontinuous on \(Q\), hence continuous; the uniform estimate
\eqref{eq:p071-exact-envelope-uniform-terminal-trace} then gives joint
continuity at each full endpoint.

It remains only to check the common node.  On a full branch, the same
monotonicity gives
\[
                         C_b(t,y)\leq C(\eta,y).
\]
The node slice is continuous by the preceding compact-polytope argument.
Along any sequence converging to the node, each subsequence either has a
further subsequence in one fixed branch, where branch-prism continuity
applies, or has a further subsequence whose full-branch labels escape every
finite set, where the last display gives the required upper bound.  Lower
semicontinuity of \(C\) gives the opposite bound.  Therefore \(C\) is
continuous at the node as well.

We have shown that \(C\in\mathcal A(D\times Q)\).  Since \(C\leq\phi\), it
is itself a member of the family in
\eqref{eq:p071-exact-envelope-definition}, and hence is its largest member.
Finally, if
\(\lVert\phi-\psi\rVert\leq\varepsilon\), then
\(C^{\mathrm{ex}}\phi-\varepsilon\) is a comparison function below
\(\psi\).  Maximality gives
\(C^{\mathrm{ex}}\phi-\varepsilon\leq C^{\mathrm{ex}}\psi\); interchanging
\(\phi\) and \(\psi\) proves
\eqref{eq:p071-exact-envelope-nonexpanding}.
\end{proof}
